
\medskip
\phantomsection\label{lem:sga3-VIB-9.2.1}%
\noindent\textbf{Lemma 9.2.1.---}
\emph{Let \(X\) be an \(S\)-scheme, and let \(R\) be an
equivalence relation on \(X\) defined by the monomorphism
\(v:R\to X\times_S X\). Suppose that \(R\) is effective.
Then:}
\begin{enumerate}
\item[\textup{(i)}]\phantomsection\label{concl:sga3-VIB-9.2.1-i}
  \emph{The morphism \(v\) is an immersion.}

\item[\textup{(ii)}]\phantomsection\label{concl:sga3-VIB-9.2.1-ii}
  \emph{For \(Y=X/R\) to be separated over \(S\), it is necessary
  that \(v\) be a closed immersion; when \(Y\) represents the
  quotient sheaf of \(X\) by \(R\) for the \textup{(fpqc)}
  topology, it is also sufficient.}\footnote{Translator's note: the 1970 text states \textup{(ii)} as an
  equivalence, without the condition on \(Y\). The proof of sufficiency uses that \(p\colon X\to Y\) is a
  covering for the \textup{(fpqc)} topology, which follows from IV, 4.4.3 when \(Y\) represents the quotient
  sheaf, but not from the effectiveness of \(R\) alone. For instance, if \(k\) is a field, the morphism
  \(p\colon X=\operatorname{Spec}k[t]\to Y=\operatorname{Spec}k[t^2,t^3]\) is the quotient of \(X\) by
  \(R=X\times_YX\) in the category of schemes, \(k[t^2,t^3]\) being the set of \(f\in k[t]\) with
  \(f\otimes1=1\otimes f\) in \(k[t]\otimes_{k[t^2,t^3]}k[t]\); but \(p\) is not a covering, since \(t^4\) is
  zero on \(X\times_YY_0\) and not on \(Y_0=\operatorname{Spec}k[t^2,t^3]/(t^3)\). The Gille--Polo re-edition
  (2011), Expos\'e~VI\textsubscript{B}, following O.~Gabber, adds the condition, which holds in the proof of
  9.2~\textup{(x)}.}
\end{enumerate}

Recall (IV, \SGARefLink{sga3:iv:definition:3:3:2}{3.3.2}) that by definition the natural morphism
\[
  R\longrightarrow X\times_Y X
\]
is an isomorphism. One deduces from this (EGA~I, 5.3.5) the
following cartesian diagram:
\begin{center}
  \phantomsection\label{diag:sga3-VIB-9.2.1-effective-square}%
  \begin{tikzcd}[row sep=large,column sep=large]
    R \arrow[r,"v"] \arrow[d] &
    X\times_SX \arrow[d,"p\times p"] \\
    Y \arrow[r,"\Delta_{Y/S}"'] &
    Y\times_SY
  \end{tikzcd}
\end{center}
Then, since \(\Delta_{Y/S}\) is an immersion (EGA~I, 5.3.9), so
is \(v\). Likewise, if \(Y\) is separated over \(S\),
\(\Delta_{Y/S}\) is a closed immersion, and hence so is \(v\).
Conversely, since \(p\) is a covering for the \textup{(fpqc)}
topology (IV, \SGARefLink{sga3:iv:proposition:4:4:3}{4.4.3}), since the property of a morphism of
being a closed immersion is of local nature for the
\textup{(fpqc)} topology (EGA~IV\(_2\), 2.7.1), and since
\(p\times_Sp\) is also a covering for the \textup{(fpqc)}
topology (IV, \SGARefLink{sga3:iv:definition:4:2:3}{4.2.3}), if \(v\) is a closed immersion, then so is
\(\Delta_{Y/S}\), and hence \(Y\) is separated over \(S\).

\medskip
\phantomsection\label{rem:sga3-VIB-9.2.2}%
\noindent\textbf{Remark 9.2.2.---}
\emph{Under the general hypotheses of
\hyperref[prop:sga3-VIB-9.2]{9.2}, if \(G'\) is flat and
locally of finite presentation over \(S\), then \(p\) is a
covering for the \textup{(fppf)} topology,}\footnote{Translator's note: the 1970 text has
``\textup{(fpqc)}'' here. To extend \textup{(vii)} and \textup{(viii)} to properties of local nature for
the \textup{(fppf)} topology, \(p\) must be a covering for \textup{(fppf)}; it is, being faithfully flat and
locally of finite presentation by 9.2~\textup{(xii)} (IV, 6.3.1~\textup{(iv)}). The Gille--Polo
re-edition (2011), Expos\'e~VI\textsubscript{B}, corrects accordingly.}
\emph{by \hyperref[concl:sga3-VIB-9.2-vii]{9.2\textup{(vii)}}.
Consequently, assertions
\hyperref[concl:sga3-VIB-9.2-vii]{9.2\textup{(vii)}} and
\hyperref[concl:sga3-VIB-9.2-viii]{9.2\textup{(viii)}} extend
to properties \(\mathscr P\) that are local for the
\textup{(fppf)} topology.}

\medskip
\phantomsection\label{rem:sga3-VIB-9.3}%
\noindent\textbf{Remark 9.3.---}
\phantomsection\label{rem:sga3-VIB-9.3-a}%
\textup{a)} The question whether a quotient \(G/G'\) is
representable is often delicate; in this seminar we prove the
representability of certain particular quotients.

In general, in order to conclude that \(G/G'\) is representable,
it is not enough to suppose that \(G\) and \(G'\) are of finite
presentation over \(S\) and that \(G'\) is flat over \(S\): even
when \(G\) and \(G'\) are smooth over \(S\), it may happen that
\(G\) has connected fibers and that \(G'\) is not closed in \(G\);
then \(G/G'\) would have to be separated
\textup{(\hyperref[thm:sga3-VIB-5.3]{5.3})},\footnote{Translator's note: the Gille--Polo
re-edition (2011), Expos\'e~VI\textsubscript{B}, refers to \hyperref[cor:sga3-VIB-5.4]{Corollary~5.4}, the
case of Theorem~5.3 that is used here: \(G\) acts on \(G/G'\), the morphism
\(G\times_S(G/G')\to(G/G')\times_S(G/G')\), \((g,x)\mapsto(gx,x)\), is surjective, and the fibers of \(G/G'\),
images of those of \(G\), are connected.} which contradicts
\hyperref[concl:sga3-VIB-9.2-x]{9.2\textup{(x)}}.

To obtain such a counterexample, take \(S\) to be the spectrum
of a discrete valuation ring and put \(G=\mathbb G_{m,S}\).
Let \(n>1\) be an integer invertible on \(S\). Then
\phantomsection\label{form:sga3-VIB-9.3-mu-n}%
\[
  \mu_n=\ker\bigl(G\xrightarrow{\ n\ }G\bigr)
\]
is a closed subgroup of \(G\), \'etale over \(S\)
(cf.\ Expos\'e~VII).\footnote{Translator's note: the Gille--Polo re-edition (2011),
Expos\'e~VI\textsubscript{B}, refers to Expos\'e~VII\(_A\), 8.4 (the case of a base field) and to
Expos\'e~VIII, \SGARefLink{sga3:viii:proposition:2:1}{2.1}. Directly: \(\mu_n=\operatorname{Spec}\mathcal O_S[T]/(T^n-1)\) is finite and of
finite presentation over \(S\), and the derivative \(nT^{n-1}\) of \(T^n-1\) is invertible on \(\mu_n\) when
\(n\) is invertible on \(S\), so that \(\mu_n\) is \'etale over \(S\).}
Let \(G'\) be the open subgroup of \(\mu_n\) obtained by
removing from \(\mu_n\) the closed subset of its closed fiber
that is complementary to the origin. Then \(G'\) is not closed
in \(G\), and hence \(G/G'\) is not representable. One can also
construct such examples in which \(G'\) is smooth with connected
fibers.

\phantomsection\label{rem:sga3-VIB-9.3-b}%
\textup{b)} It remains possible, however, that \(G/G'\) is
representable when \(G\) and \(G'\) are of finite presentation
over \(S\), and \(G'\) is flat over \(S\) and closed in \(G\)%
\begingroup
\renewcommand{\thefootnote}{\fnsymbol{footnote}}%
\footnote[1]{\label{note:sga3-VIB-9.3-star}This is too
optimistic, as M.\ Raynaud shows in his thesis
\textup{(loc.\ cit., X, 14)}. Translator's note: the thesis is M.~Raynaud,
\emph{Faisceaux amples sur les sch\'emas en groupes et les espaces homog\`enes}, Lecture Notes in Math.\
\textbf{119}, Springer, 1970; the base of its counterexample X, 14 is a regular local scheme of
dimension~\(2\), as the Gille--Polo re-edition (2011), Expos\'e~VI\textsubscript{B}, points out (see also
S.~Anantharaman, \emph{Sch\'emas en groupes, espaces homog\`enes et espaces alg\'ebriques sur une base de
dimension~1}, M\'em.\ Soc.\ Math.\ France \textbf{33} (1973), 5--79, Introduction).}%
\endgroup\hspace{0.15em}%
.
Under these hypotheses, \(G/G'\) is known to be representable
in the following particular cases:
\begin{enumerate}
\item[\textup{\(1^\circ\)---}]\phantomsection\label{case:sga3-VIB-9.3-b-1}
  \(S\) is the spectrum of an Artinian ring
  (VI\(_A\)).\footnote{Translator's note: the Gille--Polo re-edition (2011),
  Expos\'e~VI\textsubscript{B}, refers to VI\(_A\), \hyperref[thm:sga3-VIA-3.2]{3.2}, and to its VI\(_A\),
  3.3.2, which allows \(G\) not to be flat over \(S\); in this edition the latter is in
  \hyperref[thm:sga3-VIA-3.3.2]{(b) of the Addition in VI\(_A\), \S\,3}.}

\item[\textup{\(2^\circ\)---}]\phantomsection\label{case:sga3-VIB-9.3-b-2}
  \(G'\) is proper over \(S\), and \(G\) is quasi-projective
  over \(S\) (\hyperref[thm:sga3-V-7.1]{V, 7.1}).

\item[\textup{\(3^\circ\)---}]\phantomsection\label{case:sga3-VIB-9.3-b-3}
  \(S\) is regular of dimension \(1\) (Raynaud, to appear).\footnote{Translator's note: the
  Gille--Polo re-edition (2011), Expos\'e~VI\textsubscript{B}, states this case for \(S\) locally noetherian
  of dimension~\(1\), referring to S.~Anantharaman, \emph{Sch\'emas en groupes, espaces homog\`enes et espaces
  alg\'ebriques sur une base de dimension~1}, M\'em.\ Soc.\ Math.\ France \textbf{33} (1973), 5--79,
  Th\'eor\`eme~4.C: if \(S\) is locally noetherian of dimension at most \(1\), \(G\) an \(S\)-group scheme
  locally of finite type and \(H\) a closed subgroup scheme of \(G\), flat over \(S\), then the quotient
  sheaf \(G/H\) for the \textup{(fppf)} topology is a scheme.}
\end{enumerate}
